\section{认识论边界}

框架的评价分成三层：约束分类能否帮助辨认机制，数量模型能否解释并区分观察结果，规范主张能否给出公开理由。三者采用各自的证据，修订也应回到相应层面。

第一，分类框架组织问题，经过标定的数量模型才承担具体预测。预测需要固定任务、变量、基准与参数，并检验样本外表现。仅凭区域名称或未标定公式，可以列出要检查的约束，不能确定某项技术的部署年份、增长速度或临界点。

第二，物理形式进入社会技术分析，需要明确机制对应。借用剪切增稠形式时，要指出干预速率怎样引发代偿或结构变化；借用德博拉数时，要测量响应与输入的时间尺度。结构对应得到支持，模型才有解释力；对应不成立，就要更换形式。词汇相似不能代替这种检查。

第三，分类与函数形式分别接受检验。对数线性形式要与交互、阈值和其他候选形式比较，并在独立数据上检查误差；介质分类要检验是否增加了机制解释与样本外区分能力。若模型依赖反复改动变量含义才能拟合，就要重审定义。不同介质出现相似曲线，要求回到机制和任务尺度判断，曲线相似本身既不证实也不推翻整个框架。

第四，约束分析说明哪些条件影响转移代价，价值论证说明为什么某些结果、权利和责任值得被保护。本书主张让受影响者能够参与、让判断保留证据、让决定接受追责。支持这些主张，需要公开伦理理由、程序与分配安排；物理下限不替这些理由作决定。

第七章据此提出选择、分配与承担黏滞的实践主张。应用到具体决定时，先查清可行路径与成本分布，再比较结果涉及的利益和权利，最后确定授权与责任。数据支持可行性判断，规范理由支持取舍，实施记录让判断继续接受检验。

% <<<<<<<<<< parts/book-ch07.tex <<<<<<<<<<


%% ===================== 本章参考文献 =====================
\printbibliography[heading=subbibliography,title=参考文献]
